\SourceLeaf{355}%
\par\indent Ein System von Dreiecken auf dem Sphäroid, dessen Seiten sogenannte geodaetische Linien sind, wird bei einer conformen Übertragung auf die Kugelfläche durch ein analoges Dreieckssystem dargestellt, worin die Winkel, wie schon aus dem Begriffe der Conformität von selbst folgt, den entsprechenden Winkeln des erstern Systems \textit{genau} gleich sind, während die Seiten zwar nicht in mathematischer Schärfe Bögen von grössten Kreisen werden, aber doch davon nur sehr wenig abweichen. Kann man nun bewirken, dass diese Abweichungen in dem ganzen Umfange des Systems nach Massgabe der in die Berechnung zu legenden Genauigkeit wie ganz verschwindend betrachtet werden dürfen, so ist klar, dass nachdem eine Seite des sphäroidischen Systems auf die Kugelfläche übertragen ist, man ohne weiteres das ganze System wie eines von gewöhnlichen sphärischen Dreiecken berechnen darf, und nur am Schluss von den Längen und Breiten auf der Kugelfläche auf die Längen und Breiten auf dem Sphäroid zurückzugehen braucht, insofern man die Endresultate der Messung in dieser Form verlangt. Dieser Übergang wird entweder vermittelst der Formeln, welche die gewählte Übertragungsart darbietet, geschehen können, oder vermittelst einer im Voraus berechneten Hülfstafel. In den Fällen hingegen, wo jene Abweichung merklich genug wird, um eine Berücksichtigung zu verdienen, wird jeder aus den Messungen hervorgegangene Winkel vor der scharfen Berechnung auf der Kugel erst einer kleinen Reduction bedürfen, und die Arbeit wird dadurch nur unbedeutend vergrössert werden, wenn die Zahlwerthe der Reductionen sich mit Leichtigkeit berechnen lassen.

Die in der vorliegenden Abhandlung entwickelte Übertragungsart ist so beschaffen, dass die Abweichung derjenigen Curve, durch welche ein geodaetischer Bogen auf der Kugelfläche dargestellt wird, von Grösstenkreisbogen zwischen denselben Endpunkten, immer wie ganz verschwindend zu betrachten ist in der Nähe eines bestimmten Parallelkreises (Normal-Parallelkreises), welchen man nach Gefallen wählen kann, und, wenn man die ganze Rechnungsanlage von vorne her für ein bestimmtes Dreieckssystem selbst ausführt, am schicklichsten ungefähr durch die Mitte des ganzen Systems legen mag. Je weiter man sich von diesem Normal-Parallelkreise nach Norden oder Süden entfernt. desto grösser können jene Abweichungen werden. die übrigens daneben zugleich von der Grösse der Dreiecksseiten und von ihrer Lage gegen den Meridian abhängig sind; immer aber bleiben sie, selbst bei sehr beträchtlicher Entfernung von dem Normal-Pa%
